\chapter{Correlation, Transformation, and Scaling}
\label{chap:correlation-normalization}

\chapterguide
  {You should be comfortable loading data, selecting columns, using packages, and creating plots.}
  {compute and test Pearson correlation, visualize correlation matrices, and compare log transformation, standardization, and min--max scaling.}

\section{Pearson correlation with \texttt{cor()}}

Compute a Pearson coefficient for two paired vectors:

\begin{lstlisting}[style=dsr]
x <- c(1, 2, 3, 4, 5, 6, 7)
y <- c(1, 3, 6, 2, 7, 4, 5)

result <- cor(x, y, method = "pearson")
cat("Pearson correlation coefficient is:", result)
\end{lstlisting}

The coefficient summarizes their linear association from -1 to 1.

\section{Pearson correlation with \texttt{cor.test()}}

\texttt{cor.test()} adds inferential output:

\begin{lstlisting}[style=dsr]
result <- cor.test(x, y, method = "pearson")
print(result)
\end{lstlisting}

Its output includes the estimate, test statistic, p-value, and confidence interval.

\begin{pitfall}
Correlation measures association, not causation. Inspect the data and assumptions before interpreting the coefficient or p-value.
\end{pitfall}

\section{Correlation matrix with \texttt{corrplot}}

Compute and visualize the \texttt{mtcars} correlation matrix:

\begin{lstlisting}[style=dsr]
install.packages("corrplot")
install.packages("RColorBrewer")

library(corrplot)
library(RColorBrewer)

M <- cor(mtcars)
?corrplot

corrplot(M, type = "upper")
corrplot(M, type = "upper", order = "hclust")
corrplot(
  M,
  type = "upper",
  order = "hclust",
  col = brewer.pal(n = 8, name = "RdYlBu")
)
\end{lstlisting}

Hierarchical ordering places variables with similar correlation patterns together.

\section{Building a heatmap-ready correlation table}

Use the \texttt{environmental} dataset to prepare a clustered heatmap:

\begin{lstlisting}[style=dsr]
install.packages("lattice")
library(lattice)

data <- environmental
head(data)

corr_mat <- round(cor(data), 2)
head(corr_mat)
\end{lstlisting}

The matrix is reordered using a distance derived from correlation values and hierarchical clustering:

\begin{lstlisting}[style=dsr]
corr_dist <- as.dist((1 - corr_mat) / 2)
hc <- hclust(corr_dist)
corr_mat <- corr_mat[hc$order, hc$order]
\end{lstlisting}

Convert the matrix to long form for \texttt{ggplot2}:

\begin{lstlisting}[style=dsr]
install.packages("reshape2")
library(reshape2)

melted_corr_mat <- melt(corr_mat)
head(melted_corr_mat)
\end{lstlisting}


\section{Correlation heatmaps}

\subsection{Using \texttt{ggplot2}}

\begin{lstlisting}[style=dsr]
install.packages("ggplot2")
library(ggplot2)

ggplot(
  data = melted_corr_mat,
  aes(x = Var1, y = Var2, fill = value)
) +
  geom_tile() +
  geom_text(
    aes(label = value),
    color = "white",
    size = 4
  )
\end{lstlisting}

Tile color and overlaid text both encode the correlation value.

\subsection{Using \texttt{heatmaply}}

\begin{lstlisting}[style=dsr]
install.packages("heatmaply")
library(heatmaply)

heatmaply_cor(
  x = cor(data),
  xlab = "Features",
  ylab = "Features",
  k_col = 2,
  k_row = 2
)
\end{lstlisting}

\subsection{Using \texttt{ggcorrplot}}

Use the correctly named package:

\begin{lstlisting}[style=dsr]
install.packages("ggcorrplot")
library(ggcorrplot)
ggcorrplot::ggcorrplot(cor(data))
\end{lstlisting}

\begin{pitfall}
The source misspells the installation name as \texttt{ggcorplot}. The package and namespace are \texttt{ggcorrplot}.
\end{pitfall}

\section{Log transformation}

The supplied numeric vector is:

\begin{lstlisting}[style=dsr]
mydata <- c(244, 753, 596, 645, 874, 141, 639, 465, 999, 654)

scaled_data <- log(mydata)
print(scaled_data)
\end{lstlisting}

Log transformation compresses large values and can reduce right skew; it does not constrain values to a fixed range.

\section{Standardization}

\texttt{scale()} centers values at approximately zero and scales them to unit standard deviation:

\begin{lstlisting}[style=dsr]
scaled_data2 <- as.data.frame(scale(mydata))
print(scaled_data2)
\end{lstlisting}


\section{Min--max scaling}

Min--max scaling maps values to a fixed range, normally 0 to 1:

\begin{lstlisting}[style=dsr]
install.packages("caret")
library(caret)

minmax <- preProcess(
  as.data.frame(mydata),
  method = c("range")
)

scaled_data3 <- predict(
  minmax,
  as.data.frame(mydata)
)
print(scaled_data3)
\end{lstlisting}

\section{Comparing summaries}

Compare the raw and transformed summaries:

\begin{lstlisting}[style=dsr]
summary(mydata)
summary(scaled_data)
summary(scaled_data2)
summary(scaled_data3)
\end{lstlisting}

\begin{pitfall}
The source calls an undefined object, \texttt{scaled\_data1}. The corrected comparison uses \texttt{scaled\_data}.
\end{pitfall}

\begin{dsfigure}
\begin{tikzpicture}[
  box/.style={draw=DSBlue,fill=DSPaleBlue,rounded corners=2mm,text width=3.35cm,minimum height=1.1cm,align=center,font=\small\bfseries\color{DSNavy}},
  arr/.style={-{Stealth[length=2mm]},thick,draw=DSTeal}
]
\node[box] (raw) {Raw numeric data};
\node[box,right=1.0cm of raw] (log) {Log transformation};
\node[box,below=0.8cm of log] (std) {Standard scaling};
\node[box,above=0.8cm of log] (mm) {Min--max scaling};
\draw[arr] (raw) -- (log);
\draw[arr] (raw.east) -- ++(0.55,0) |- (std.west);
\draw[arr] (raw.east) -- ++(0.55,0) |- (mm.west);
\end{tikzpicture}
\end{dsfigure}

Compare each method with the raw data because the resulting scales have different interpretations.

\begin{quickcheck}
\begin{enumerate}
  \item What is the difference in output intent between \texttt{cor()} and \texttt{cor.test()} in Lab 9a?
  \item Which dataset is used for the first \texttt{corrplot()} example?
  \item Why is \texttt{melt()} applied to the reordered correlation matrix before the \texttt{ggplot2} heatmap?
  \item Which function performs the standard scaling in the handout?
  \item Which object-name mismatch appears in the final summary comparison?
\end{enumerate}
\end{quickcheck}

\begin{chapterrecap}
\begin{itemize}
  \item \texttt{cor()} estimates association; \texttt{cor.test()} adds inferential statistics.
  \item Clustered heatmaps group variables with similar correlation patterns.
  \item Log transformation, standardization, and min--max scaling serve different purposes.
\end{itemize}
\end{chapterrecap}
